\Chapter{2}

\Verse{1}
{
    \zA{却说封肃听见公差传唤，忙出来陪笑启问，那些人只嚷：“快请出甄爷来。” 封肃忙陪笑道：“小人姓封，并不姓甄。 只有当日小婿姓甄，今已出家一二年了，
        不知可是问他？” 那些公人道：“我们也不知什么‘真’‘假’，既是你的女婿， 就带了你去面？ 太爷便了。” 大家把封肃推拥而去，封家各各惊慌，不知何事。
        至 二更时分，封肃方回来，众人忙问端的。}
}
{
    \eA{To continue. Feng Su, upon hearing the shouts of the public messengers, came
        out in a flurry and forcing a smile, he asked them to explain (their
        errand); but all these people did was to continue bawling out: "Be quick,
        and ask Mr. Chen to come out." "My surname is Feng," said Feng Su, as he
        promptly forced himself to smile;}
}
{
}

\Verse{2}
{
    \zA{――“原来新任太爷姓贾名化，本湖州人 氏，曾与女婿旧交，因在我家门首看见娇杏丫头买线，只说女婿移住此间，所以来 传。
        我将缘故回明，那太爷感伤叹息了一回； 又问外孙女儿，我说看灯丢了。 太爷 说：‘不妨，待我差人去，务必找寻回来。’ 说了一回话，临走又送我二两银子。”
        甄家娘子听了，不觉感伤。 一夜无话。 次日，早有雨村遣人送了两封银子、四匹锦缎，答谢甄家娘子； 又一封密书与 封肃，托他向甄家娘子要那娇杏作二房。}
}
{
    \eA{"It is'nt Chen at all: I had once a son-in-law whose surname was Chen, but
        he has left home, it is now already a year or two back. Is it perchance
        about him that you are inquiring?" To which the public servants remarked:
        "We know nothing about Chen or Chia (true or false); but as he is your
        son-in-law, we'll take you at once along with us to make verbal answer to
        our master and have done with it."}
}
{
}

\Verse{3}
{
    \zA{封肃喜得眉开眼笑，巴不得去奉承太爷， 便在女儿前一力撺掇。 当夜用一乘小轿，便把娇杏送进衙内去了。 雨村欢喜自不必
        言，又封百金赠与封肃，又送甄家娘子许多礼物，令其且自过活，以待访寻女儿下 落。 却说娇杏那丫头便是当年回顾雨村的，因偶然一看便弄出这段奇缘，也是意想
        不到之事。 谁知他命运两济，不承望自到雨村身边，只一年便生一子，又半载雨村 嫡配忽染疾下世，雨村便将他扶作正室夫人。 正是： 偶因一回顾，便为人上人。}
}
{
    \eA{And forthwith the whole bevy of public servants hustled Feng Su on, as they
        went on their way back; while every one in the Feng family was seized with
        consternation, and could not imagine what it was all about. It was no
        earlier than the second watch, when Feng Su returned home; and they, one and
        all, pressed him with questions as to what had happened.}
}
{
}

\Verse{4}
{
    \zA{原来雨村因那年士隐赠银之后，他于十六日便起身赴京。 大比之期，十分得意， 中了进士，选入外班，今已升了本县太爷。 虽才干优长，未免贪酷，且恃才侮上，
        那同寅皆侧目而视。 不上一年，便被上司参了一本，说他貌似有才，性实狡猾，又 题了一两件徇庇蠹役、交结乡绅之事，龙颜大怒，即命革职。 部文一到，本府各官
        无不喜悦。 那雨村虽十分惭恨，面上却全无一点怨色，仍是嘻笑自若；}
}
{
    \eA{"The fact is," he explained, "the newly-appointed Magistrate, whose surname
        is Chia, whose name is Huo and who is a native of Hu-chow, has been on
        intimate terms, in years gone by, with our son-in-law; that at the sight of
        the girl Chiao Hsing, standing at the door, in the act of buying thread, he
        concluded that he must have shifted his quarters over here, and hence it was
        that his messengers came to fetch him.}
}
{
}

\Verse{5}
{
    \zA{交代过了公 事，将历年所积的宦囊，并家属人等，送至原籍安顿妥当了，却自己担风袖月，游 览天下胜迹。 那日偶又游至维扬地方，闻得今年盐政点的是林如海。
        这林如海姓林名海，表 字如海，乃是前科的探花，今已升兰台寺大夫，本贯姑苏人氏，今钦点为巡盐御史， 到任未久。
        原来这林如海之祖也曾袭过列侯的，今到如海，业经五世。 起初只袭三 世，因当今隆恩盛德，额外加恩，至如海之父又袭了一代，到了如海便从科第出身。
        虽系世禄之家，却是书香之族。}
}
{
    \eA{I gave him a clear account of the various circumstances (of his
        misfortunes), and the Magistrate was for a time much distressed and
        expressed his regret. He then went on to make inquiries about my
        grand-daughter, and I explained that she had been lost, while looking at the
        illuminations.}
}
{
}

\Verse{6}
{
    \zA{只可惜这林家支庶不盛，人丁有限，虽有几门，却 与如海俱是堂族，没甚亲支嫡派的。 今如海年已五十，只有一个三岁之子，又于去
        岁亡了，虽有几房姬妾，奈命中无子，亦无可如何之事。 只嫡妻贾氏生得一女，乳 名黛玉，年方五岁，夫妻爱之如掌上明珠。 见他生得聪明俊秀，也欲使他识几个字，
        不过假充养子，聊解膝下荒凉之叹。 且说贾雨村在旅店偶感风寒，愈后又因盘费不继，正欲得一个居停之所以为息 肩之地。}
}
{
    \eA{'No matter,' put in the Magistrate, 'I will by and by order my men to make
        search, and I feel certain that they will find her and bring her back.' Then
        ensued a short conversation, after which I was about to go, when he
        presented me with the sum of two taels." The mistress of the Chen family
        (Mrs. Chen Shih-yin) could not but feel very much affected by what she
        heard, and the whole evening she uttered not a word.}
}
{
}

\Verse{7}
{
    \zA{偶遇两个旧友认得新盐政，知他正要请一西席教训女儿，遂将雨村荐进衙 门去。 这女学生年纪幼小，身体又弱，工课不限多寡，其馀不过两个伴读丫鬟，故
        雨村十分省力，正好养病。 看看又是一载有馀，不料女学生之母贾氏夫人一病而亡。 女学生奉侍汤药，守丧尽礼，过于哀痛，素本怯弱，因此旧病复发，有好些时不曾
        上学。 雨村闲居无聊，每当风日晴和，饭后便出来闲步。 这一日偶至郊外，意欲赏鉴那村野风光。}
}
{
    \eA{The next day, at an early hour, Yü-ts'un sent some of his men to bring over
        to Chen's wife presents, consisting of two packets of silver, and four
        pieces of brocaded silk, as a token of gratitude, and to Feng Su also a
        confidential letter, requesting him to ask of Mrs. Chen her maid Chiao Hsing
        to become his second wife.}
}
{
}

\Verse{8}
{
    \zA{信步至一山环水漩、茂林修竹之处， 隐隐有座庙宇，门巷倾颓，墙垣剥落。 有额题曰“智通寺”。 门旁又有一副旧破的 对联云： 身后有馀忘缩手，
        眼前无路想回头。 雨村看了，因想道：“这两句文虽甚浅，其意则深。 也曾游过些名山大刹，倒不曾
        见过这话头，其中想必有个翻过筋斗来的也未可知，何不进去一访。” 走入看时， 只有一个龙钟老僧在那里煮粥。 雨村见了，却不在意； 及至问他两句话，那老僧既
        聋且昏，又齿落舌钝，所答非所问。}
}
{
    \eA{Feng Su was so intensely delighted that his eyebrows expanded, his eyes
        smiled, and he felt eager to toady to the Magistrate (by presenting the girl
        to him). He hastened to employ all his persuasive powers with his daughter
        (to further his purpose), and on the same evening he forthwith escorted
        Chiao Hsing in a small chair to the Yamên. The joy experienced by Yü-ts'un
        need not be dilated upon.}
}
{
}

\Verse{9}
{
    \zA{雨村不耐烦，仍退出来，意欲到那村肆中沽饮 三杯，以助野趣。 于是移步行来。 刚入肆门，只见座上吃酒之客有一人起身大笑，
        接了出来，口内说：“奇遇，奇遇！” 雨村忙看时，此人是都中古董行中贸易姓冷 号子兴的，旧日在都相识。 雨村最赞这冷子兴是个有作为大本领的人，这子兴又借
        雨村斯文之名，故二人最相投契。 雨村忙亦笑问：“老兄何日到此？ 弟竟不知。 今 日偶遇，真奇缘也。”
        子兴道：“去年岁底到家，今因还要入都，从此顺路找个敝 友说一句话。}
}
{
    \eA{He also presented Feng Su with a packet containing one hundred ounces of
        gold; and sent numerous valuable presents to Mrs. Chen, enjoining her "to
        live cheerfully in the anticipation of finding out the whereabouts of her
        daughter."}
}
{
}

\Verse{10}
{
    \zA{承他的情，留我多住两日。 我也无甚紧事，且盘桓两日，待月半时也 就起身了。 今日敝友有事，我因闲走到此，不期这样巧遇！” 一面说一面让雨村同
        席坐了，另整上酒肴来。 二人闲谈慢饮，叙些别后之事。 雨村因问：“近日都中可有新闻没有？” 子兴
        道：“倒没有什么新闻，倒是老先生的贵同宗家出了一件小小的异事。” 雨村笑道： “弟族中无人在都，何谈及此？” 子兴笑道：“你们同姓，岂非一族？”
        雨村问： “是谁家？”}
}
{
    \eA{It must be explained, however, that the maid Chi'ao Hsing was the very
        person, who, a few years ago, had looked round at Yü-ts'un and who, by one
        simple, unpremeditated glance, evolved, in fact, this extraordinary destiny
        which was indeed an event beyond conception.}
}
{
}

\Verse{11}
{
    \zA{子兴笑道：“荣国贾府中，可也不玷辱老先生的门楣了！” 雨村道： “原来是他家。 若论起来，寒族人丁却自不少，东汉贾复以来，支派繁盛，各省皆
        有，谁能逐细考查？ 若论荣国一支，却是同谱。 但他那等荣耀，我们不便去认他， 故越发生疏了。” 子兴叹道：“老先生休这样说。
        如今的这荣、宁两府，也都萧索 了，不比先时的光景！” 雨村道：“当日宁荣两宅人口也极多，如何便萧索了呢？” 子兴道：“正是，说来也话长。”}
}
{
    \eA{Who would ever have foreseen that fate and fortune would both have so
        favoured her that she should, contrary to all anticipation, give birth to a
        son, after living with Yü-ts'un barely a year, that in addition to this,
        after the lapse of another half year, Yü-ts'un's wife should have contracted
        a sudden illness and departed this life, and that Yü-ts'un should have at
        once raised her to the rank of first wife.}
}
{
}

\Verse{12}
{
    \zA{雨村道：“去岁我到金陵时，因欲游览六朝遗迹， 那日进了石头城，从他宅门前经过。 街东是宁国府，街西是荣国府，二宅相连，竟 将大半条街占了。
        大门外虽冷落无人，隔着围墙一望，里面厅殿楼阁也还都峥嵘轩 峻，就是后边一带花园里，树木山石，也都还有葱蔚洇润之气，那里像个衰败之 家？”
        子兴笑道：“亏你是进士出身，原来不通。 古人有言：‘百足之虫，死而不 僵。’ 如今虽说不似先年那样兴盛，较之平常仕宦人家，到底气象不同。}
}
{
    \eA{Her destiny is adequately expressed by the lines:  Through but one single,
        casual look  Soon an exalted place she took. The fact is that after Yü-ts'un
        had been presented with the money by Shih-yin, he promptly started on the
        16th day for the capital, and at the triennial great tripos, his wishes were
        gratified to the full.}
}
{
}

\Verse{13}
{
    \zA{如今人口 日多，事务日盛，主仆上下都是安富尊荣，运筹谋画的竟无一个，那日用排场，又 不能将就省俭。 如今外面的架子虽没很倒，内囊却也尽上来了。
        这也是小事。 更有 一件大事：谁知这样钟鸣鼎食的人家儿，如今养的儿孙，竟一代不如一代了！” 雨村听说，也道：“这样诗礼之家，岂有不善教育之理？
        别门不知，只说这宁 荣两宅，是最教子有方的，何至如此？” 子兴叹道：“正说的是这两门呢。 等我告 诉你：当日宁国公是一母同胞弟兄两个。}
}
{
    \eA{Having successfully carried off his degree of graduate of the third rank,
        his name was put by selection on the list for provincial appointments. By
        this time, he had been raised to the rank of Magistrate in this district;
        but, in spite of the excellence and sufficiency of his accomplishments and
        abilities, he could not escape being ambitious and overbearing.}
}
{
}

\Verse{14}
{
    \zA{宁公居长，生了两个儿子。 宁公死后，长 子贾代化袭了官，也养了两个儿子：长子名贾敷，八九岁上死了，只剩了一个次子
        贾敬，袭了官，如今一味好道，只爱烧丹炼汞，别事一概不管。 幸而早年留下一个 儿子，名唤贾珍，因他父亲一心想作神仙，把官倒让他袭了。 他父亲又不肯住在家
        里，只在都中城外和那些道士们胡羼。 这位珍爷也生了一个儿子，今年才十六岁， 名叫贾蓉。 如今敬老爷不管事了，这珍爷那里干正事？}
}
{
    \eA{He failed besides, confident as he was in his own merits, in respect toward
        his superiors, with the result that these officials looked upon him
        scornfully with the corner of the eye.}
}
{
}

\Verse{15}
{
    \zA{只一味高乐不了，把那宁国 府竟翻过来了，也没有敢来管他的人。 再说荣府你听：方才所说异事就出在这里。
        自荣公死后，长子贾代善袭了官，娶的是金陵世家史侯的小姐为妻。 生了两个儿子， 长名贾赦，次名贾政。 如今代善早已去世，太夫人尚在。
        长子贾赦袭了官，为人却 也中平，也不管理家事； 惟有次子贾政，自幼酷喜读书，为人端方正直。 祖父钟爱，
        原要他从科甲出身，不料代善临终遗本一上，皇上怜念先臣，即叫长子袭了官；}
}
{
    \eA{A year had hardly elapsed, when he was readily denounced in a memorial to
        the Throne by the High Provincial authorities, who represented that he was
        of a haughty disposition, that he had taken upon himself to introduce
        innovations in the rites and ceremonies, that overtly, while he endeavoured
        to enjoy the reputation of probity and uprightness, he, secretly, combined
        the nature of the tiger and wolf;}
}
{
}

\Verse{16}
{
    \zA{又 问还有几个儿子，立刻引见，又将这政老爷赐了个额外主事职衔，叫他入部习学， 如今现已升了员外郎。
        这政老爷的夫人王氏，头胎生的公子名叫贾珠，十四岁进学， 后来娶了妻、生了子，不到二十岁，一病就死了。 第二胎生了一位小姐，生在大年 初一，就奇了。
        不想隔了十几年，又生了一位公子，说来更奇：一落胞胎，嘴里便 衔下一块五彩晶莹的玉来，还有许多字迹。 你道是新闻不是？”
        雨村笑道：“果然奇异，只怕这人的来历不小。”}
}
{
    \eA{with the consequence that he had been the cause of much trouble in the
        district, and that he had made life intolerable for the people, \&c. \&c.
        The Dragon countenance of the Emperor was considerably incensed. His Majesty
        lost no time in issuing commands, in reply to the Memorial, that he should
        be deprived of his official status.}
}
{
}

\Verse{17}
{
    \zA{子兴冷笑道：“万人都这样 说，因而他祖母爱如珍宝。 那周岁时，政老爷试他将来的志向，便将世上所有的东 西摆了无数叫他抓。
        谁知他一概不取，伸手只把些脂粉钗环抓来玩弄，那政老爷便 不喜欢，说将来不过酒色之徒，因此不甚爱惜。 独那太君还是命根子一般。 ――说
        来又奇：如今长了十来岁，虽然淘气异常，但聪明乖觉，百个不及他一个； 说起孩 子话来也奇，他说：‘女儿是水做的骨肉，男子是泥做的骨肉。 我见了女儿便清爽，
        见了男子便觉浊臭逼人。’}
}
{
    \eA{On the arrival of the despatch from the Board, great was the joy felt by
        every officer, without exception, of the prefecture in which he had held
        office. Yü-ts'un, though at heart intensely mortified and incensed, betrayed
        not the least outward symptom of annoyance, but still preserved, as of old,
        a smiling and cheerful countenance.}
}
{
}

\Verse{18}
{
    \zA{你道好笑不好笑？ 将来色鬼无疑了！” 雨村罕然厉色道：“非也!可惜你们不知道这人的来历，大约政老前辈也错以 淫魔色鬼看待了。
        若非多读书识事，加以致知格物之功、悟道参玄之力者，不能知 也。” 子兴见他说得这样重大，忙请教其故。 雨村道：“天地生人，除大仁大恶， 馀者皆无大异。
        若大仁者则应运而生，大恶者则应劫而生，运生世治，劫生世危。 尧、舜、禹、汤、文、武、周、召、孔、孟、董、韩、周、程、朱、张，皆应运而 生者；}
}
{
    \eA{He handed over charge of all official business and removed the savings which
        he had accumulated during the several years he had been in office, his
        family and all his chattels to his original home; where, after having put
        everything in proper order, he himself travelled (carried the winds and
        sleeved the moon) far and wide, visiting every relic of note in the whole
        Empire.}
}
{
}

\Verse{19}
{
    \zA{蚩尤、共工、桀、纣、始皇、王莽、曹操、桓温、安禄山、秦桧等，皆应劫 而生者。 大仁者修治天下，大恶者扰乱天下。 清明灵秀，天地之正气，仁者之所秉 也；
        残忍乖僻，天地之邪气，恶者之所秉也。 今当祚永运隆之日，太平无为之世， 清明灵秀之气所秉者，上自朝廷，下至草野，比比皆是。 所馀之秀气漫无所归，遂
        为甘露、为和风，洽然溉及四海。 彼残忍乖邪之气，不能荡溢于光天化日之下，遂 凝结充塞于深沟大壑之中。}
}
{
    \eA{As luck would have it, on a certain day while making a second journey
        through the Wei Yang district, he heard the news that the Salt Commissioner
        appointed this year was Lin Ju-hai. This Lin Ju-hai's family name was Lin,
        his name Hai and his style Ju-hai. He had obtained the third place in the
        previous triennial examination, and had, by this time, already risen to the
        rank of Director of the Court of Censors.}
}
{
}

\Verse{20}
{
    \zA{偶因风荡，或被云摧，略有摇动感发之意，一丝半缕误 而逸出者，值灵秀之气适过，正不容邪，邪复妒正，两不相下； 如风水雷电地中既
        遇，既不能消，又不能让，必致搏击掀发。 既然发泄，那邪气亦必赋之于人。 假使 或男或女偶秉此气而生者，上则不能为仁人为君子，下亦不能为大凶大恶。 置之千
        万人之中，其聪俊灵秀之气，则在千万人之上； 其乖僻邪谬不近人情之态，又在千 万人之下。 若生于公侯富贵之家，则为情痴情种。 若生于诗书清贫之族，则为逸士
        高人。}
}
{
    \eA{He was a native of Kú Su. He had been recently named by Imperial appointment
        a Censor attached to the Salt Inspectorate, and had arrived at his post only
        a short while back. In fact, the ancestors of Lin Ju-hai had, from years
        back, successively inherited the title of Marquis, which rank, by its
        present descent to Ju-hai, had already been enjoyed by five generations.}
}
{
}

\Verse{21}
{
    \zA{纵然生于薄祚寒门，甚至为奇优，为名娼，亦断不至为走卒健仆，甘遭庸夫 驱制。 如前之许由、陶潜、阮籍、嵇康、刘伶、王谢二族、顾虎头、陈后主、唐明
        皇、宋徽宗、刘庭芝、温飞卿、米南宫、石曼卿、柳耆卿、秦少游，近日倪云林、 唐伯虎、祝枝山，再如李龟年、黄幡绰、敬新磨、卓文君、红拂、薛涛、崔莺、朝
        云之流，此皆易地则同之人也。” 子兴道：“依你说，‘成则公侯败则贼’了？” 雨村道：“正是这意。}
}
{
    \eA{When first conferred, the hereditary right to the title had been limited to
        three generations; but of late years, by an act of magnanimous favour and
        generous beneficence, extraordinary bounty had been superadded; and on the
        arrival of the succession to the father of Ju-hai, the right had been
        extended to another degree.}
}
{
}

\Verse{22}
{
    \zA{你还不 知，我自革职以来，这两年遍游各省，也曾遇见两个异样孩子，所以方才你一说这 宝玉，我就猜着了八九也是这一派人物。
        不用远说，只这金陵城内钦差金陵省体仁 院总裁甄家，你可知道？” 子兴道：“谁人不知!这甄府就是贾府老亲，他们两家 来往极亲热的。
        就是我也和他家往来非止一日了。” 雨村笑道：“去岁我在金陵， 也曾有人荐我到甄府处馆。 我进去看其光景，谁知他家那等荣贵，却是个富而好礼
        之家，倒是个难得之馆。}
}
{
    \eA{It had now descended to Ju-hai, who had, besides this title of nobility,
        begun his career as a successful graduate. But though his family had been
        through uninterrupted ages the recipient of imperial bounties, his kindred
        had all been anyhow men of culture. The only misfortune had been that the
        several branches of the Lin family had not been prolific, so that the
        numbers of its members continued limited;}
}
{
}

\Verse{23}
{
    \zA{但是这个学生虽是启蒙，却比一个举业的还劳神，说起来 更可笑，他说：‘必得两个女儿陪着我读书，我方能认得字，心上也明白，不然我 心里自己糊涂。’
        又常对着跟他的小厮们说：‘这女儿两个字极尊贵极清净的，比 那瑞兽珍禽、奇花异草更觉希罕尊贵呢，你们这种浊口臭舌万万不可唐突了这两个
        字，要紧，要紧!但凡要说的时节，必用净水香茶漱了口方可； 设若失错，便要凿 牙穿眼的。’ 其暴虐顽劣，种种异常；}
}
{
    \eA{and though there existed several households, they were all however to Ju-hai
        no closer relatives than first cousins. Neither were there any connections
        of the same lineage, or of the same parentage. Ju-hai was at this date past
        forty; and had only had a son, who had died the previous year, in the third
        year of his age. Though he had several handmaids, he had not had the good
        fortune of having another son;}
}
{
}

\Verse{24}
{
    \zA{只放了学进去，见了那些女儿们，其温厚和 平、聪敏文雅，竟变了一个样子。 因此他令尊也曾下死笞楚过几次，竟不能改。 每
        打的吃疼不过时，他便‘姐姐’‘妹妹’的乱叫起来。 后来听得里面女儿们拿他取 笑：‘因何打急了只管叫姐妹作什么？ 莫不叫姐妹们去讨情讨饶？ 你岂不愧些！’
        他 回答的最妙，他说：‘急痛之时，只叫姐姐妹妹字样，或可解疼也未可知，因叫了 一声，果觉疼得好些。 遂得了秘法，每疼痛之极，便连叫姐妹起来了。’
        你说可笑 不可笑？}
}
{
    \eA{but this was too a matter that could not be remedied. By his wife, née Chia,
        he had a daughter, to whom the infant name of Tai Yü was given. She was, at
        this time, in her fifth year. Upon her the parents doated as much as if she
        were a brilliant pearl in the palm of their hand.}
}
{
}

\Verse{25}
{
    \zA{为他祖母溺爱不明，每因孙辱师责子，我所以辞了馆出来的。 这等子弟必 不能守祖父基业、从师友规劝的。 只可惜他家几个好姊妹都是少有的！”
        子兴道：“便是贾府中现在三个也不错。 政老爷的长女名元春，因贤孝才德， 选入宫作女史去了。 二小姐乃是赦老爷姨娘所出，名迎春。 三小姐政老爷庶出，名
        探春。 四小姐乃宁府珍爷的胞妹，名惜春。 因史老夫人极爱孙女，都跟在祖母这边， 一处读书，听得个个不错。”}
}
{
    \eA{Seeing that she was endowed with natural gifts of intelligence and good
        looks, they also felt solicitous to bestow upon her a certain knowledge of
        books, with no other purpose than that of satisfying, by this illusory way,
        their wishes of having a son to nurture and of dispelling the anguish felt
        by them, on account of the desolation and void in their family circle (round
        their knees). But to proceed.}
}
{
}

\Verse{26}
{
    \zA{雨村道：“更妙在甄家风俗，女儿之名亦皆从男子之 名，不似别人家里另外用这些‘春’‘红’‘香’‘玉’等艳字。 何得贾府亦落此 俗套？” 子兴道：“不然。
        只因现今大小姐是正月初一所生，故名‘元春’，馀者 都从了‘春’字； 上一排的却也是从弟兄而来的。 现有对证：目今你贵东家林公的
        夫人，即荣府中赦、政二公的胞妹，在家时名字唤贾敏。 不信时你回去细访可知。” 雨村拍手笑道：“是极。}
}
{
    \eA{Yü-ts'un, while sojourning at an inn, was unexpectedly laid up with a
        violent chill. Finding on his recovery, that his funds were not sufficient
        to pay his expenses, he was thinking of looking out for some house where he
        could find a resting place when he suddenly came across two friends
        acquainted with the new Salt Commissioner.}
}
{
}

\Verse{27}
{
    \zA{我这女学生名叫黛玉，他读书凡‘敏’字他皆念作‘密’ 字，写字遇着‘敏’字亦减一二笔。 我心中每每疑惑，今听你说，是为此无疑矣。
        怪道我这女学生言语举止另是一样，不与凡女子相同。 度其母不凡，故生此女，今 知为荣府之外孙，又不足罕矣!可惜上月其母竟亡故了。” 子兴叹道：“老姊妹三
        个，这是极小的，又没了!长一辈的姊妹一个也没了。 只看这小一辈的，将来的东 床何如呢。” 雨村道：“正是。}
}
{
    \eA{Knowing that this official was desirous to find a tutor to instruct his
        daughter, they lost no time in recommending Yü-ts'un, who moved into the
        Yamên. His female pupil was youthful in years and delicate in physique, so
        that her lessons were irregular.}
}
{
}

\Verse{28}
{
    \zA{方才说政公已有一个衔玉之子，又有长子所遗弱孙，这赦老 竟无一个不成？” 子兴道：“政公既有玉儿之后，其妾又生了一个，倒不知其好歹。
        只眼前现有二子一孙，却不知将来何如。 若问那赦老爷，也有一子，名叫贾琏，今 已二十多岁了，亲上做亲，娶的是政老爷夫人王氏内侄女，今已娶了四五年。 这位
        琏爷身上现捐了个同知，也是不喜正务的，于世路上好机变，言谈去得，所以目今 现在乃叔政老爷家住，帮着料理家务。}
}
{
    \eA{Besides herself, there were only two waiting girls, who remained in
        attendance during the hours of study, so that Yü-ts'un was spared
        considerable trouble and had a suitable opportunity to attend to the
        improvement of his health. In a twinkle, another year and more slipped by,
        and when least expected, the mother of his ward, née Chia, was carried away
        after a short illness.}
}
{
}

\Verse{29}
{
    \zA{谁知自娶了这位奶奶之后，倒上下无人不称 颂他的夫人，琏爷倒退了一舍之地：模样又极标致，言谈又爽利，心机又极深细， 竟是个男人万不及一的。”
        雨村听了笑道：“可知我言不谬了。 你我方才所说的这 几个人，只怕都是那正邪两赋而来一路之人，未可知也。”
        子兴道：“正也罢，邪也罢，只顾算别人家的账，你也吃杯酒才好。” 雨村道： “只顾说话，就多吃了几杯。” 子兴笑道：“说着别人家的闲话，正好下酒，即多
        吃几杯何妨。”}
}
{
    \eA{His pupil (during her mother's sickness) was dutiful in her attendance, and
        prepared the medicines for her use. (And after her death,) she went into the
        deepest mourning prescribed by the rites, and gave way to such excess of
        grief that, naturally delicate as she was, her old complaint, on this
        account, broke out anew.}
}
{
}

\Verse{30}
{
    \zA{雨村向窗外看道：“天也晚了，仔细关了城，我们慢慢进城再谈， 未为不可。” 于是二人起身，算还酒钱。 方欲走时，忽听得后面有人叫道：“雨村
        兄恭喜了!特来报个喜信的。” 雨村忙回头看时，――要知是谁，且听下回分解。 (本章完)}
}
{
    \eA{Being unable for a considerable time to prosecute her studies, Yü-ts'un
        lived at leisure and had no duties to attend to. Whenever therefore the wind
        was genial and the sun mild, he was wont to stroll at random, after he had
        done with his meals.}
}
{
}

\Verse{31}
{
    \zA{}
}
{
    \eA{On this particular day, he, by some accident, extended his walk beyond the
        suburbs, and desirous to contemplate the nature of the rustic scenery, he,
        with listless step, came up to a spot encircled by hills and streaming
        pools, by luxuriant clumps of trees and thick groves of bamboos. Nestling in
        the dense foliage stood a temple. The doors and courts were in ruins. The
        walls, inner and outer, in disrepair. An inscription on a tablet testified
        that this was the temple of Spiritual Perception. On the sides of the door
        was also a pair of old and dilapidated scrolls with the following
        enigmatical verses. Behind ample there is, yet to retract the hand, the mind
        heeds not,  until. Before the mortal vision lies no path, when comes to turn
        the will.}
}
{
}

\Verse{32}
{
    \zA{}
}
{
    \eA{"These two sentences," Yü-ts'un pondered after perusal, "although simple in
        language, are profound in signification. I have previous to this visited
        many a spacious temple, located on hills of note, but never have I beheld an
        inscription referring to anything of the kind. The meaning contained in
        these words must, I feel certain, owe their origin to the experiences of
        some person or other; but there's no saying. But why should I not go in and
        inquire for myself?" Upon walking in, he at a glance caught sight of no one
        else, but of a very aged bonze, of unkempt appearance, cooking his rice.}
}
{
}

\Verse{33}
{
    \zA{}
}
{
    \eA{When Yü-ts'un perceived that he paid no notice, he went up to him and asked
        him one or two questions, but as the old priest was dull of hearing and a
        dotard, and as he had lost his teeth, and his tongue was blunt, he made most
        irrelevant replies. Yü-ts'un lost all patience with him, and withdrew again
        from the compound with the intention of going as far as the village public
        house to have a drink or two, so as to enhance the enjoyment of the rustic
        scenery. With easy stride, he accordingly walked up to the place.}
}
{
}

\Verse{34}
{
    \zA{}
}
{
    \eA{Scarcely had he passed the threshold of the public house, when he perceived
        some one or other among the visitors who had been sitting sipping their wine
        on the divan, jump up and come up to greet him, with a face beaming with
        laughter. "What a strange meeting! What a strange meeting!" he exclaimed
        aloud. Yü-ts'un speedily looked at him, (and remembered) that this person
        had, in past days, carried on business in a curio establishment in the
        capital, and that his surname was Leng and his style Tzu-hsing. A mutual
        friendship had existed between them during their sojourn, in days of yore,
        in the capital;}
}
{
}

\Verse{35}
{
    \zA{}
}
{
    \eA{and as Yü-ts'un had entertained the highest opinion of Leng Tzu-hsing, as
        being a man of action and of great abilities, while this Leng Tzu-hsing, on
        the other hand, borrowed of the reputation of refinement enjoyed by
        Yü-ts'un, the two had consequently all along lived in perfect harmony and
        companionship. "When did you get here?" Yü-ts'un eagerly inquired also
        smilingly. "I wasn't in the least aware of your arrival. This unexpected
        meeting is positively a strange piece of good fortune." "I went home,"
        Tzu-hsing replied, "about the close of last year, but now as I am again
        bound to the capital, I passed through here on my way to look up a friend of
        mine and talk some matters over.}
}
{
}

\Verse{36}
{
    \zA{}
}
{
    \eA{He had the kindness to press me to stay with him for a couple of days
        longer, and as I after all have no urgent business to attend to, I am
        tarrying a few days, but purpose starting about the middle of the moon. My
        friend is busy to-day, so I roamed listlessly as far as here, never dreaming
        of such a fortunate meeting." While speaking, he made Yü-ts'un sit down at
        the same table, and ordered a fresh supply of wine and eatables; and as the
        two friends chatted of one thing and another, they slowly sipped their wine.
        The conversation ran on what had occurred after the separation, and Yü-ts'un
        inquired, "Is there any news of any kind in the capital?" "There's nothing
        new whatever," answered Tzu-hsing.}
}
{
}

\Verse{37}
{
    \zA{}
}
{
    \eA{"There is one thing however: in the family of one of your worthy kinsmen, of
        the same name as yourself, a trifling, but yet remarkable, occurrence has
        taken place." "None of my kindred reside in the capital," rejoined Yü-ts'un
        with a smile. "To what can you be alluding?" "How can it be that you people
        who have the same surname do not belong to one clan?" remarked Tzu-hsing,
        sarcastically. "In whose family?" inquired Yü-ts'un. "The Chia family,"
        replied Tzu-hsing smiling, "whose quarters are in the Jung Kuo Mansion, does
        not after all reflect discredit upon the lintel of your door, my venerable
        friend." "What!" exclaimed Yü-ts'un, "did this affair take place in that
        family? Were we to begin reckoning, we would find the members of my clan to
        be anything but limited in number.}
}
{
}

\Verse{38}
{
    \zA{}
}
{
    \eA{Since the time of our ancestor Chia Fu, who lived while the Eastern Han
        dynasty occupied the Throne, the branches of our family have been numerous
        and flourishing; they are now to be found in every single province, and who
        could, with any accuracy, ascertain their whereabouts? As regards the
        Jung-kuo branch in particular, their names are in fact inscribed on the same
        register as our own, but rich and exalted as they are, we have never
        presumed to claim them as our relatives, so that we have become more and
        more estranged."}
}
{
}

\Verse{39}
{
    \zA{}
}
{
    \eA{"Don't make any such assertions," Tzu-hsing remarked with a sigh, "the
        present two mansions of Jung and Ning have both alike also suffered
        reverses, and they cannot come up to their state of days of yore." "Up to
        this day, these two households of Ning and of Jung," Yü-ts'un suggested,
        "still maintain a very large retinue of people, and how can it be that they
        have met with reverses?" "To explain this would be indeed a long story,"
        said Leng Tzu-hsing. "Last year," continued Yü-ts'un, "I arrived at Chin
        Ling, as I entertained a wish to visit the remains of interest of the six
        dynasties, and as I on that day entered the walled town of Shih T'ou, I
        passed by the entrance of that old residence. On the east side of the
        street, stood the Ning Kuo mansion; on the west the Jung Kuo mansion;}
}
{
}

\Verse{40}
{
    \zA{}
}
{
    \eA{and these two, adjoining each other as they do, cover in fact well-nigh half
        of the whole length of the street. Outside the front gate everything was, it
        is true, lonely and deserted; but at a glance into the interior over the
        enclosing wall, I perceived that the halls, pavilions, two-storied
        structures and porches presented still a majestic and lofty appearance. Even
        the flower garden, which extends over the whole area of the back grounds,
        with its trees and rockeries, also possessed to that day an air of
        luxuriance and freshness, which betrayed no signs of a ruined or decrepid
        establishment."}
}
{
}

\Verse{41}
{
    \zA{}
}
{
    \eA{"You have had the good fortune of starting in life as a graduate," explained
        Tzu-tsing as he smiled, "and yet are not aware of the saying uttered by some
        one of old: that a centipede even when dead does not lie stiff. (These
        families) may, according to your version, not be up to the prosperity of
        former years, but, compared with the family of an ordinary official, their
        condition anyhow presents a difference. Of late the number of the inmates
        has, day by day, been on the increase; their affairs have become daily more
        numerous; of masters and servants, high and low, who live in ease and
        respectability very many there are; but of those who exercise any
        forethought, or make any provision, there is not even one.}
}
{
}

\Verse{42}
{
    \zA{}
}
{
    \eA{In their daily wants, their extravagances, and their expenditure, they are
        also unable to adapt themselves to circumstances and practise economy; (so
        that though) the present external framework may not have suffered any
        considerable collapse, their purses have anyhow begun to feel an exhausting
        process! But this is a mere trifle. There is another more serious matter.
        Would any one ever believe that in such families of official status, in a
        clan of education and culture, the sons and grandsons of the present age
        would after all be each (succeeding) generation below the standard of the
        former?"}
}
{
}

\Verse{43}
{
    \zA{}
}
{
    \eA{Yü-ts'un, having listened to these remarks, observed: "How ever can it be
        possible that families of such education and refinement can observe any
        system of training and nurture which is not excellent? Concerning the other
        branches, I am not in a position to say anything; but restricting myself to
        the two mansions of Jung and Ning, they are those in which, above all
        others, the education of their children is methodical." "I was just now
        alluding to none other than these two establishments," Tzu-hsing observed
        with a sigh; "but let me tell you all. In days of yore, the duke of Ning Kuo
        and the duke of Jung Kuo were two uterine brothers. The Ning duke was the
        elder; he had four sons.}
}
{
}

\Verse{44}
{
    \zA{}
}
{
    \eA{After the death of the duke of Ning Kuo, his eldest son, Chia Tai-hua, came
        into the title. He also had two sons; but the eldest, whose name was Hu,
        died at the age of eight or nine; and the only survivor, the second son,
        Chia Ching, inherited the title. His whole mind is at this time set upon
        Taoist doctrines; his sole delight is to burn the pill and refine the dual
        powers; while every other thought finds no place in his mind. Happily, he
        had, at an early age, left a son, Chia Chen, behind in the lay world, and
        his father, engrossed as his whole heart was with the idea of attaining
        spiritual life, ceded the succession of the official title to him.}
}
{
}

\Verse{45}
{
    \zA{}
}
{
    \eA{His parent is, besides, not willing to return to the original family seat,
        but lives outside the walls of the capital, foolishly hobnobbing with all
        the Taoist priests. This Mr. Chen had also a son, Chia Jung, who is, at this
        period, just in his sixteenth year. Mr. Ching gives at present no attention
        to anything at all, so that Mr. Chen naturally devotes no time to his
        studies, but being bent upon nought else but incessant high pleasure, he has
        subversed the order of things in the Ning Kuo mansion, and yet no one can
        summon the courage to come and hold him in check. But I'll now tell you
        about the Jung mansion for your edification.}
}
{
}

\Verse{46}
{
    \zA{}
}
{
    \eA{The strange occurrence, to which I alluded just now, came about in this
        manner. After the demise of the Jung duke, the eldest son, Chia Tai-shan,
        inherited the rank. He took to himself as wife, the daughter of Marquis
        Shih, a noble family of Chin Ling, by whom he had two sons; the elder being
        Chia She, the younger Chia Cheng. This Tai Shan is now dead long ago; but
        his wife is still alive, and the elder son, Chia She, succeeded to the
        degree. He is a man of amiable and genial disposition, but he likewise gives
        no thought to the direction of any domestic concern. The second son Chia
        Cheng displayed, from his early childhood, a great liking for books, and
        grew up to be correct and upright in character.}
}
{
}

\Verse{47}
{
    \zA{}
}
{
    \eA{His grandfather doated upon him, and would have had him start in life
        through the arena of public examinations, but, when least expected,
        Tai-shan, being on the point of death, bequeathed a petition, which was laid
        before the Emperor. His Majesty, out of regard for his former minister,
        issued immediate commands that the elder son should inherit the estate, and
        further inquired how many sons there were besides him, all of whom he at
        once expressed a wish to be introduced in his imperial presence.}
}
{
}

\Verse{48}
{
    \zA{}
}
{
    \eA{His Majesty, moreover, displayed exceptional favour, and conferred upon Mr.
        Cheng the brevet rank of second class Assistant Secretary (of a Board), and
        commanded him to enter the Board to acquire the necessary experience. He has
        already now been promoted to the office of second class Secretary. This Mr.
        Cheng's wife, nèe Wang, first gave birth to a son called Chia Chu, who
        became a Licentiate in his fourteenth year. At barely twenty, he married,
        but fell ill and died soon after the birth of a son. Her (Mrs. Cheng's)
        second child was a daughter, who came into the world, by a strange
        coincidence, on the first day of the year.}
}
{
}

\Verse{49}
{
    \zA{}
}
{
    \eA{She had an unexpected (pleasure) in the birth, the succeeding year, of
        another son, who, still more remarkable to say, had, at the time of his
        birth, a piece of variegated and crystal-like brilliant jade in his mouth,
        on which were yet visible the outlines of several characters. Now, tell me,
        was not this a novel and strange occurrence? eh?" "Strange indeed!"
        exclaimed Yü-ts'un with a smile; "but I presume the coming experiences of
        this being will not be mean." Tzu-hsing gave a faint smile. "One and all,"
        he remarked, "entertain the same idea. Hence it is that his mother doats
        upon him like upon a precious jewel.}
}
{
}

\Verse{50}
{
    \zA{}
}
{
    \eA{On the day of his first birthday, Mr. Cheng readily entertained a wish to
        put the bent of his inclinations to the test, and placed before the child
        all kinds of things, without number, for him to grasp from. Contrary to
        every expectation, he scorned every other object, and, stretching forth his
        hand, he simply took hold of rouge, powder and a few hair-pins, with which
        he began to play. Mr. Cheng experienced at once displeasure, as he
        maintained that this youth would, by and bye, grow up into a sybarite,
        devoted to wine and women, and for this reason it is, that he soon began to
        feel not much attachment for him. But his grandmother is the one who, in
        spite of everything, prizes him like the breath of her own life. The very
        mention of what happened is even strange!}
}
{
}

\Verse{51}
{
    \zA{}
}
{
    \eA{He is now grown up to be seven or eight years old, and, although
        exceptionally wilful, in intelligence and precocity, however, not one in a
        hundred could come up to him! And as for the utterances of this child, they
        are no less remarkable. The bones and flesh of woman, he argues, are made of
        water, while those of man of mud. 'Women to my eyes are pure and pleasing,'
        he says, 'while at the sight of man, I readily feel how corrupt, foul and
        repelling they are!' Now tell me, are not these words ridiculous? There can
        be no doubt whatever that he will by and bye turn out to be a licentious
        roué." Yü-ts'un, whose countenance suddenly assumed a stern air, promptly
        interrupted the conversation. "It doesn't quite follow," he suggested.}
}
{
}

\Verse{52}
{
    \zA{}
}
{
    \eA{"You people don't, I regret to say, understand the destiny of this child.
        The fact is that even the old Hanlin scholar Mr. Cheng was erroneously
        looked upon as a loose rake and dissolute debauchee! But unless a person,
        through much study of books and knowledge of letters, so increases (in lore)
        as to attain the talent of discerning the nature of things, and the vigour
        of mind to fathom the Taoist reason as well as to comprehend the first
        principle, he is not in a position to form any judgment."}
}
{
}

\Verse{53}
{
    \zA{}
}
{
    \eA{Tzu-hsing upon perceiving the weighty import of what he propounded, "Please
        explain," he asked hastily, "the drift (of your argument)." To which
        Yü-ts'un responded: "Of the human beings created by the operation of heaven
        and earth, if we exclude those who are gifted with extreme benevolence and
        extreme viciousness, the rest, for the most part, present no striking
        diversity. If they be extremely benevolent, they fall in, at the time of
        their birth, with an era of propitious fortune; while those extremely
        vicious correspond, at the time of their existence, with an era of calamity.
        When those who coexist with propitious fortune come into life, the world is
        in order; when those who coexist with unpropitious fortune come into life,
        the world is in danger.}
}
{
}

\Verse{54}
{
    \zA{}
}
{
    \eA{Yao, Shun, Yü, Ch'eng T'ang, Wen Wang, Wu Wang, Chou Kung, Chao Kung,
        Confucius, Mencius, T'ung Hu, Han Hsin, Chou Tzu, Ch'eng Tzu, Chu Tzu and
        Chang Tzu were ordained to see light in an auspicious era. Whereas Ch'i Yu,
        Kung Kung, Chieh Wang, Chou Wang, Shih Huang, Wang Mang, Tsao Ts'ao, Wen
        Wen, An Hu-shan, Ch'in Kuei and others were one and all destined to come
        into the world during a calamitous age. Those endowed with extreme
        benevolence set the world in order; those possessed of extreme maliciousness
        turn the world into disorder. Purity, intelligence, spirituality and
        subtlety constitute the vital spirit of right which pervades heaven and
        earth, and the persons gifted with benevolence are its natural fruit.}
}
{
}

\Verse{55}
{
    \zA{}
}
{
    \eA{Malignity and perversity constitute the spirit of evil, which permeates
        heaven and earth, and malicious persons are affected by its influence. The
        days of perpetual happiness and eminent good fortune, and the era of perfect
        peace and tranquility, which now prevail, are the offspring of the pure,
        intelligent, divine and subtle spirit which ascends above, to the very
        Emperor, and below reaches the rustic and uncultured classes. Every one is
        without exception under its influence.}
}
{
}

\Verse{56}
{
    \zA{}
}
{
    \eA{The superfluity of the subtle spirit expands far and wide, and finding
        nowhere to betake itself to, becomes, in due course, transformed into dew,
        or gentle breeze; and, by a process of diffusion, it pervades the whole
        world. "The spirit of malignity and perversity, unable to expand under the
        brilliant sky and transmuting sun, eventually coagulates, pervades and stops
        up the deep gutters and extensive caverns; and when of a sudden the wind
        agitates it or it be impelled by the clouds, and any slight disposition, on
        its part, supervenes to set itself in motion, or to break its bounds, and so
        little as even the minutest fraction does unexpectedly find an outlet, and
        happens to come across any spirit of perception and subtlety which may be at
        the time passing by, the spirit of right does not yield to the spirit of
        evil, and the spirit of evil is again envious of the spirit of right, so
        that the two do not harmonize.}
}
{
}

\Verse{57}
{
    \zA{}
}
{
    \eA{Just like wind, water, thunder and lightning, which, when they meet in the
        bowels of the earth, must necessarily, as they are both to dissolve and are
        likewise unable to yield, clash and explode to the end that they may at
        length exhaust themselves. Hence it is that these spirits have also forcibly
        to diffuse themselves into the human race to find an outlet, so that they
        may then completely disperse, with the result that men and women are
        suddenly imbued with these spirits and spring into existence. At best,
        (these human beings) cannot be generated into philanthropists or perfect
        men; at worst, they cannot also embody extreme perversity or extreme
        wickedness.}
}
{
}

\Verse{58}
{
    \zA{}
}
{
    \eA{Yet placed among one million beings, the spirit of intelligence, refinement,
        perception and subtlety will be above these one million beings; while, on
        the other hand, the perverse, depraved and inhuman embodiment will likewise
        be below the million of men. Born in a noble and wealthy family, these men
        will be a salacious, lustful lot; born of literary, virtuous or poor
        parentage, they will turn out retired scholars or men of mark; though they
        may by some accident be born in a destitute and poverty-stricken home, they
        cannot possibly, in fact, ever sink so low as to become runners or menials,
        or contentedly brook to be of the common herd or to be driven and curbed
        like a horse in harness.}
}
{
}

\Verse{59}
{
    \zA{}
}
{
    \eA{They will become, for a certainty, either actors of note or courtesans of
        notoriety; as instanced in former years by Hsü Yu, T'ao Ch'ien, Yuan Chi,
        Chi Kang, Liu Ling, the two families of Wang and Hsieh, Ku Hu-t'ou, Ch'en
        Hou-chu, T'ang Ming-huang, Sung Hui-tsung, Liu T'ing-chih, Wen Fei-ching,
        Mei Nan-kung, Shih Man-ch'ing, Lui C'hih-ch'ing and Chin Shao-yu, and
        exemplified now-a-days by Ni Yün-lin, T'ang Po-hu, Chu Chih-shan, and also
        by Li Kuei-men, Huang P'an-cho, Ching Hsin-mo, Cho Wen-chün;}
}
{
}

\Verse{60}
{
    \zA{}
}
{
    \eA{and the women Hung Fu, Hsieh T'ao, Ch'ü Ying, Ch'ao Yün and others; all of
        whom were and are of the same stamp, though placed in different scenes of
        action." "From what you say," observed Tzu-hsing, "success makes (a man) a
        duke or a marquis; ruin, a thief!" "Quite so; that's just my idea!" replied
        Yü-ts'un; "I've not as yet let you know that after my degradation from
        office, I spent the last couple of years in travelling for pleasure all over
        each province, and that I also myself came across two extraordinary youths.
        This is why, when a short while back you alluded to this Pao-yü, I at once
        conjectured, with a good deal of certainty, that he must be a human being of
        the same stamp.}
}
{
}

\Verse{61}
{
    \zA{}
}
{
    \eA{There's no need for me to speak of any farther than the walled city of Chin
        Ling. This Mr. Chen was, by imperial appointment, named Principal of the
        Government Public College of the Chin Ling province. Do you perhaps know
        him?" "Who doesn't know him?" remarked Tzu-hsing. "This Chen family is an
        old connection of the Chia family. These two families were on terms of great
        intimacy, and I myself likewise enjoyed the pleasure of their friendship for
        many a day." "Last year, when at Chin Ling," Yü-ts'un continued with a
        smile, "some one recommended me as resident tutor to the school in the Chen
        mansion; and when I moved into it I saw for myself the state of things. Who
        would ever think that that household was grand and luxurious to such a
        degree!}
}
{
}

\Verse{62}
{
    \zA{}
}
{
    \eA{But they are an affluent family, and withal full of propriety, so that a
        school like this was of course not one easy to obtain. The pupil, however,
        was, it is true, a young tyro, but far more troublesome to teach than a
        candidate for the examination of graduate of the second degree. Were I to
        enter into details, you would indeed have a laugh. 'I must needs,' he
        explained, 'have the company of two girls in my studies to enable me to read
        at all, and to keep likewise my brain clear. Otherwise, if left to myself,
        my head gets all in a muddle.'}
}
{
}

\Verse{63}
{
    \zA{}
}
{
    \eA{Time after time, he further expounded to his young attendants, how extremely
        honourable and extremely pure were the two words representing woman, that
        they are more valuable and precious than the auspicious animal, the
        felicitous bird, rare flowers and uncommon plants. 'You may not' (he was
        wont to say), 'on any account heedlessly utter them, you set of foul mouths
        and filthy tongues! these two words are of the utmost import! Whenever you
        have occasion to allude to them, you must, before you can do so with
        impunity, take pure water and scented tea and rinse your mouths. In the
        event of any slip of the tongue, I shall at once have your teeth extracted,
        and your eyes gouged out.' His obstinacy and waywardness are, in every
        respect, out of the common.}
}
{
}

\Verse{64}
{
    \zA{}
}
{
    \eA{After he was allowed to leave school, and to return home, he became, at the
        sight of the young ladies, so tractable, gentle, sharp, and polite,
        transformed, in fact, like one of them. And though, for this reason, his
        father has punished him on more than one occasion, by giving him a sound
        thrashing, such as brought him to the verge of death, he cannot however
        change. Whenever he was being beaten, and could no more endure the pain, he
        was wont to promptly break forth in promiscuous loud shouts, 'Girls! girls!'
        The young ladies, who heard him from the inner chambers, subsequently made
        fun of him. 'Why,' they said, 'when you are being thrashed, and you are in
        pain, your only thought is to bawl out girls! Is it perchance that you
        expect us young ladies to go and intercede for you?}
}
{
}

\Verse{65}
{
    \zA{}
}
{
    \eA{How is that you have no sense of shame?' To their taunts he gave a most
        plausible explanation. 'Once,' he replied, 'when in the agony of pain, I
        gave vent to shouting girls, in the hope, perchance, I did not then know, of
        its being able to alleviate the soreness. After I had, with this purpose,
        given one cry, I really felt the pain considerably better; and now that I
        have obtained this secret spell, I have recourse, at once, when I am in the
        height of anguish, to shouts of girls, one shout after another. Now what do
        you say to this? Isn't this absurd, eh?"}
}
{
}

\Verse{66}
{
    \zA{}
}
{
    \eA{"The grandmother is so infatuated by her extreme tenderness for this youth,
        that, time after time, she has, on her grandson's account, found fault with
        the tutor, and called her son to task, with the result that I resigned my
        post and took my leave. A youth, with a disposition such as his, cannot
        assuredly either perpetuate intact the estate of his father and grandfather,
        or follow the injunctions of teacher or advice of friends. The pity is,
        however, that there are, in that family, several excellent female cousins,
        the like of all of whom it would be difficult to discover." "Quite so!"
        remarked Tzu-hsing;}
}
{
}

\Verse{67}
{
    \zA{}
}
{
    \eA{"there are now three young ladies in the Chia family who are simply
        perfection itself. The eldest is a daughter of Mr. Cheng, Yuan Ch'un by
        name, who, on account of her excellence, filial piety, talents, and virtue,
        has been selected as a governess in the palace. The second is the daughter
        of Mr. She's handmaid, and is called Ying Ch'un; the third is T'an Ch'un,
        the child of Mr. Cheng's handmaid; while the fourth is the uterine sister of
        Mr. Chen of the Ning Mansion. Her name is Hsi Ch'un. As dowager lady Shih is
        so fondly attached to her granddaughters, they come, for the most part, over
        to their grandmother's place to prosecute their studies together, and each
        one of these girls is, I hear, without a fault."}
}
{
}

\Verse{68}
{
    \zA{}
}
{
    \eA{"More admirable," observed Yü-ts'un, "is the régime (adhered to) in the Chen
        family, where the names of the female children have all been selected from
        the list of male names, and are unlike all those out-of-the-way names, such
        as Spring Blossom, Scented Gem, and the like flowery terms in vogue in other
        families. But how is it that the Chia family have likewise fallen into this
        common practice?" "Not so!" ventured Tzu-h'sing. "It is simply because the
        eldest daughter was born on the first of the first moon, that the name of
        Yuan Ch'un was given to her; while with the rest this character Ch'un
        (spring) was then followed. The names of the senior generation are, in like
        manner, adopted from those of their brothers; and there is at present an
        instance in support of this.}
}
{
}

\Verse{69}
{
    \zA{}
}
{
    \eA{The wife of your present worthy master, Mr. Lin, is the uterine sister of
        Mr. Chia. She and Mr. Chia Cheng, and she went, while at home, under the
        name of Chia Min. Should you question the truth of what I say, you are at
        liberty, on your return, to make minute inquiries and you'll be convinced."
        Yü-ts'un clapped his hands and said smiling, "It's so, I know! for this
        female pupil of mine, whose name is Tai-yü, invariably pronounces the
        character \_min\_ as \_mi\_, whenever she comes across it in the course of
        her reading; while, in writing, when she comes to the character 'min,' she
        likewise reduces the strokes by one, sometimes by two.}
}
{
}

\Verse{70}
{
    \zA{}
}
{
    \eA{Often have I speculated in my mind (as to the cause), but the remarks I've
        heard you mention, convince me, without doubt, that it is no other reason
        (than that of reverence to her mother's name). Strange enough, this pupil of
        mine is unique in her speech and deportment, and in no way like any ordinary
        young lady. But considering that her mother was no commonplace woman
        herself, it is natural that she should have given birth to such a child.
        Besides, knowing, as I do now, that she is the granddaughter of the Jung
        family, it is no matter of surprise to me that she is what she is. Poor
        girl, her mother, after all, died in the course of the last month."
        Tzu-hsing heaved a sigh.}
}
{
}

\Verse{71}
{
    \zA{}
}
{
    \eA{"Of three elderly sisters," he explained, "this one was the youngest, and
        she too is gone! Of the sisters of the senior generation not one even
        survives! But now we'll see what the husbands of this younger generation
        will be like by and bye!" "Yes," replied Yü-ts'un. "But some while back you
        mentioned that Mr. Cheng has had a son, born with a piece of jade in his
        mouth, and that he has besides a tender-aged grandson left by his eldest
        son; but is it likely that this Mr. She has not, himself, as yet, had any
        male issue?" "After Mr. Cheng had this son with the jade," Tzu-hsing added,
        "his handmaid gave birth to another son, who whether he be good or bad, I
        don't at all know.}
}
{
}

\Verse{72}
{
    \zA{}
}
{
    \eA{At all events, he has by his side two sons and a grandson, but what these
        will grow up to be by and bye, I cannot tell. As regards Mr. Chia She, he
        too has had two sons; the second of whom, Chia Lien, is by this time about
        twenty. He took to wife a relative of his, a niece of Mr. Cheng's wife, a
        Miss Wang, and has now been married for the last two years. This Mr. Lien
        has lately obtained by purchase the rank of sub-prefect. He too takes little
        pleasure in books, but as far as worldly affairs go, he is so versatile and
        glib of tongue, that he has recently taken up his quarters with his uncle
        Mr. Cheng, to whom he gives a helping hand in the management of domestic
        matters.}
}
{
}

\Verse{73}
{
    \zA{}
}
{
    \eA{Who would have thought it, however, ever since his marriage with his worthy
        wife, not a single person, whether high or low, has there been who has not
        looked up to her with regard: with the result that Mr. Lien himself has, in
        fact, had to take a back seat (\_lit\_. withdrew 35 li). In looks, she is
        also so extremely beautiful, in speech so extremely quick and fluent, in
        ingenuity so deep and astute, that even a man could, in no way, come up to
        her mark." After hearing these remarks Yü-ts'un smiled.}
}
{
}

\Verse{74}
{
    \zA{}
}
{
    \eA{"You now perceive," he said, "that my argument is no fallacy, and that the
        several persons about whom you and I have just been talking are, we may
        presume, human beings, who, one and all, have been generated by the spirit
        of right, and the spirit of evil, and come to life by the same royal road;
        but of course there's no saying." "Enough," cried Tzu-hsing, "of right and
        enough of evil; we've been doing nothing but settling other people's
        accounts; come now, have another glass, and you'll be the better for it!"
        "While bent upon talking," Yü-ts'un explained, "I've had more glasses than
        is good for me." "Speaking of irrelevant matters about other people,"
        Tzu-hsing rejoined complacently, "is quite the thing to help us swallow our
        wine; so come now;}
}
{
}

\Verse{75}
{
    \zA{}
}
{
    \eA{what harm will happen, if we do have a few glasses more." Yü-ts'un thereupon
        looked out of the window. "The day is also far advanced," he remarked, "and
        if we don't take care, the gates will be closing; let us leisurely enter the
        city, and as we go along, there will be nothing to prevent us from
        continuing our chat." Forthwith the two friends rose from their seats,
        settled and paid their wine bill, and were just going, when they
        unexpectedly heard some one from behind say with a loud voice: "Accept my
        congratulations, Brother Yü-ts'un; I've now come, with the express purpose
        of giving you the welcome news!" Yü-ts'un lost no time in turning his head
        round to look at the speaker.}
}
{
}

\Verse{76}
{
    \zA{}
}
{
    \eA{But reader, if you wish to learn who the man was, listen to the details
        given in the following chapter.}
}
{
}
